\documentclass[journal,10pt]{IEEEtran}
\setlength{\textwidth}{182mm}
\setlength{\columnsep}{4.2mm}
\usepackage{amsmath,amssymb,newtxmath}
\usepackage[no-math]{fontspec}
\setmainfont{Times New Roman}
\usepackage{xeCJK}
\setCJKmainfont[AutoFakeBold=2]{SimSun}
\usepackage{graphicx,booktabs,array,hyperref,placeins}
\hypersetup{hidelinks,unicode=true}
\raggedbottom
\setlength{\emergencystretch}{2em}
\renewcommand{\arraystretch}{1.15}
\usepackage{seqsplit}
\begin{document}
\title{Litz 线的分层验证、共端复电流分配与有限长度纵向 PEEC · 03 结果}
\author{缪雨轩\thanks{淮南师范学院}}

\maketitle
\noindent 合并稿对应章节 / Chapter in merged manuscript: \href{paper_zh-CN_ieee.pdf#chapter-03}{03}\par
\FloatBarrier
\section*{03 结果}
\hypertarget{chapter-03}{}
\pdfbookmark[0]{03 结果}{bookmark-chapter-03}
\subsection*{单圆线解析与独立二维校准}
\hypertarget{03.analytic.heading}{}
\pdfbookmark[1]{单圆线解析与独立二维校准}{bookmark-03.analytic.heading}
长直圆线内部解采用修正 Bessel 函数 \href{paper_zh-CN_ieee.pdf#E03}{(6)}。200 kHz 下的趋肤深度 δ=0.147771655 mm，相对 1 mm 半径已形成明显表层电流集中。解析值为 Rdc=5.488101486 mΩ/m、Rac=20.016203628 mΩ/m、Rac/Rdc=3.647199980、损耗 8.006481451 W/m。这里的内部阻抗不含由回流确定的外部电感。径向相位保留原复相量，并通过相邻半径的 360° 连续提升显示；显示解缠不修改数据。


独立二维圆盘扩散使用三档 353、1389、5573 节点，并按总复电流归一化。最细 Rac 解析相对误差为 1.79233×10⁻⁶，功率积分差为 3.06800×10⁻⁶。结果支持该孤立内场问题的离散与归一化。二维场与有限长度三维端区不具有相同边界，不能把这个误差直接用于后者。解析截面图由径向解析幅值线性插值并轴对称映射。


\subsection*{640 股二维结果及其边界}
\hypertarget{03.coax_results.heading}{}
\pdfbookmark[1]{640 股二维结果及其边界}{bookmark-03.coax_results.heading}
640 股直束的冻结三档网格、回流与外域检查共有 23 项通过。最细网格增量的 Rac 变化为 0.011902\%，回流距离变化为 0.000495\%，外域变化为 0.000270\%，均低于本路线 1\% 门槛。它们说明指定二维共端模型在已检查扰动范围内稳定，不能代表真实绕制束的有限端区已经收敛。束内 Rac 与回流损耗分别报告，避免把回流中的热损耗算成导体自身改善。


逐股幅值与相位展示了共同端压下的电流分配，虚线 It/N 仅为参照，不是求解约束。节点电流积分与公共端条件的一致性，是均流讨论的必要前提。对这种二维直束，把单股直径减小并不等价于沿程位置已经换位。它的主要用途是验证多股场路与受控边界，作为后续换位模型的截面证据；最终结构排序仍需要保留沿程耦合、实际端接与相关数值不确定性的模型或实验。


\subsection*{纵向 PEEC 敏感性与结果}
\hypertarget{03.peec_results.heading}{}
\pdfbookmark[1]{纵向 PEEC 敏感性与结果}{bookmark-03.peec_results.heading}
1264 自由度的 0.16、0.32、0.64 m 三档 Rac 分别为 19.899215、19.958642、19.988477 mΩ/轴向m。两个长度增量的变化均低于 1\%；5024 自由度的 0.32\(\to\)0.64 m 变化也低于门槛。互项 q4/6/8、自项 10/8\(\to\)12/10 以及当前实际截面的粗细离散比较均达到专项规定的有限检查目标。试验逐次改变一个因素；未运行的 1.28 m 档保留为未运行，不能填入外推值或零误差。


末档 5024 自由度、Lz=0.64 m、互项 q4、自项 12/10 的 Rac=20.012738128 mΩ/轴向m、Rac/Rdc=3.644719868，损耗为 8.005095251 W/轴向m；与本档等面积 Bessel 参照的差为 0.044582\%。截面颜色图按保存的细丝电流和方形单元恢复。该细丝电流偏差量描述离散导体内部的电流集中，不是 Litz 物理绝缘股线的均流指标。


长度敏感性与解析参照分别回答不同问题。前者检验有限长度核在所选区间内的变化，后者核验电阻是否与可比较的长直内场尺度一致。二者接近并不消除边界或基函数缺口。特别是末档的很小积分变化只能支持当前积分实现，对真实端区、曲线股线的横向电流以及端接再分布没有新增证据。论文因而把这一结果称为纵向模型内的有限验收。


\subsection*{CPU/GPU 等价性与资源边界}
\hypertarget{03.compute.heading}{}
\pdfbookmark[1]{CPU/GPU 等价性与资源边界}{bookmark-03.compute.heading}
CPU/GPU 对照固定同一纵向模型和离散设置。5024 自由度的 L 矩阵相对差为 7.48149×10⁻¹⁶，复电流相对差为 7.78082×10⁻¹⁴，低于 10⁻¹⁰ 门槛；1264 对照亦通过。这支持已实现 CPU/GPU 算法在所比较矩阵与端口解上的等价。它不赋予模型额外物理内容，也不能用不同规模或不同积分设置的用时声称加速比。


9856 自由度历史结果只恢复了部分电流等数据，缺少完整矩阵、公共电压或全场的项目不升级为完整验证。20108 及更大规模保持已有资源边界证据和退出记录。资源曲线只采用可归属的进程树样本；短时峰值可能遗漏，混入其他 PID 的记录排除。计算能力边界与物理失败分开解释：缺少资源不是通过，运行完成也不是物理验收通过。


640 股束内 Rac=15.793003197 mΩ/轴向m，束内损耗 6.317201279 W/轴向m，回流损耗 2.196205516 W/轴向m，实际铜截面 5.018478 mm²。共同端轴向电场为 0.411856245+j(4.064969103) V/m。归一化复电流 RMS 偏差为 1.854214，各股电流的复数和恢复 20 A RMS 总输运。


\hypertarget{E03}{}
\begin{equation}
\begin{gathered}
\delta=(\pi f\mu\sigma)^{-1/2} \\
\kappa=(1+j)/\delta \\
J_z(r)=\frac{I_t\kappa I_0(\kappa r)}{2\pi a I_1(\kappa a)} \\
z_{\mathrm{int}}=\frac{\kappa I_0(\kappa a)}{2\pi a\sigma I_1(\kappa a)}
\end{gathered}
\tag{6}\label{eq:E03}
\end{equation}
\begin{figure*}[!htbp]
\centering
\setcounter{figure}{2}
\refstepcounter{figure}\label{fig:q1_profile}
\hypertarget{fig-q1_profile}{}
\includegraphics[width=182mm]{../../../composites/zh-CN/q1_profile.pdf}
\par\vspace{4pt}{\small\centering 图 3. 单圆线径向电流幅值与相位。\par}
\vspace{3pt}{\small\raggedright 铜导体直径 2 mm，200 kHz、20 A RMS。 (a) RMS 幅值；(b) 连续显示的相位。\par}
\end{figure*}
\begin{figure}[!htbp]
\centering
\setcounter{figure}{3}
\refstepcounter{figure}\label{fig:q1_cross_section}
\hypertarget{fig-q1_cross_section}{}
\includegraphics[width=88.9mm]{../../../panels/zh-CN/q1_cross_section.pdf}
\par\vspace{4pt}{\small\centering 图 4. 单圆线解析截面幅值。\par}
\vspace{3pt}{\small\raggedright 铜导体直径 2 mm，200 kHz、20 A RMS。 色标显示轴向 RMS 电流密度。\par}
\end{figure}
\begin{figure}[!htbp]
\centering
\setcounter{figure}{4}
\refstepcounter{figure}\label{fig:q1_fem_error}
\hypertarget{fig-q1_fem_error}{}
\includegraphics[width=88.9mm]{../../../panels/zh-CN/q1_fem_error.pdf}
\par\vspace{4pt}{\small\centering 图 5. 单圆线二维校准误差。\par}
\vspace{3pt}{\small\raggedright 铜导体直径 2 mm，200 kHz、20 A RMS。 三档网格与独立功率积分对应同一内场问题。\par}
\end{figure}
\begin{figure*}[!htbp]
\centering
\setcounter{figure}{5}
\refstepcounter{figure}\label{fig:coax640_currents}
\hypertarget{fig-coax640_currents}{}
\includegraphics[width=182mm]{../../../composites/zh-CN/coax640_currents.pdf}
\par\vspace{4pt}{\small\centering 图 6. 640 股共端复电流分配。\par}
\vspace{3pt}{\small\raggedright (a) 幅值；(b) 相位，按物理股线身份排列。\par}
\end{figure*}
\begin{figure*}[!htbp]
\centering
\setcounter{figure}{6}
\refstepcounter{figure}\label{fig:coax640_fields}
\hypertarget{fig-coax640_fields}{}
\includegraphics[width=182mm]{../../../composites/zh-CN/coax640_fields.pdf}
\par\vspace{4pt}{\small\centering 图 7. 640 股二维截面电流密度。\par}
\vspace{3pt}{\small\raggedright (a) RMS 幅值；(b) 相位，材料标记识别各股。\par}
\end{figure*}
\begin{figure}[!htbp]
\centering
\setcounter{figure}{7}
\refstepcounter{figure}\label{fig:coax640_mesh}
\hypertarget{fig-coax640_mesh}{}
\includegraphics[width=88.9mm]{../../../panels/zh-CN/coax640_mesh.pdf}
\par\vspace{4pt}{\small\centering 图 8. 640 股网格与边界敏感性。\par}
\vspace{3pt}{\small\raggedright 各试验按网格、回流和外域分别改变一个因素。\par}
\end{figure}
\begin{figure}[!htbp]
\centering
\setcounter{figure}{8}
\refstepcounter{figure}\label{fig:peec_cross_section}
\hypertarget{fig-peec_cross_section}{}
\includegraphics[width=88.9mm]{../../../panels/zh-CN/peec_cross_section.pdf}
\par\vspace{4pt}{\small\centering 图 9. 纵向 PEEC 截面电流。\par}
\vspace{3pt}{\small\raggedright 方形细丝单元的复电流恢复 RMS 幅值。\par}
\end{figure}
\begin{figure}[!htbp]
\centering
\setcounter{figure}{9}
\refstepcounter{figure}\label{fig:peec_length}
\hypertarget{fig-peec_length}{}
\includegraphics[width=88.9mm]{../../../panels/zh-CN/peec_length.pdf}
\par\vspace{4pt}{\small\centering 图 10. 纵向 PEEC 长度敏感性。\par}
\vspace{3pt}{\small\raggedright 同一截面按保存的轴向长度比较单位长度电阻。\par}
\end{figure}
\begin{figure}[!htbp]
\centering
\setcounter{figure}{10}
\refstepcounter{figure}\label{fig:peec_sensitivity}
\hypertarget{fig-peec_sensitivity}{}
\includegraphics[width=88.9mm]{../../../panels/zh-CN/peec_sensitivity.pdf}
\par\vspace{4pt}{\small\centering 图 11. 纵向 PEEC 积分敏感性。\par}
\vspace{3pt}{\small\raggedright 互项阶数与自项精度各自进行单因素对照。\par}
\end{figure}
\begin{figure}[!htbp]
\centering
\setcounter{figure}{11}
\refstepcounter{figure}\label{fig:gpu_equivalence}
\hypertarget{fig-gpu_equivalence}{}
\includegraphics[width=88.9mm]{../../../panels/zh-CN/gpu_equivalence.pdf}
\par\vspace{4pt}{\small\centering 图 12. CPU/GPU 同模型数值对照。\par}
\vspace{3pt}{\small\raggedright 矩阵与复电流相对差使用同一离散和积分设置。\par}
\end{figure}
\begin{figure}[!htbp]
\centering
\setcounter{figure}{12}
\refstepcounter{figure}\label{fig:resource_memory}
\hypertarget{fig-resource_memory}{}
\includegraphics[width=88.9mm]{../../../panels/zh-CN/resource_memory.pdf}
\par\vspace{4pt}{\small\centering 图 13. 可归属进程树资源曲线。\par}
\vspace{3pt}{\small\raggedright 曲线使用归档的时间及内存样本。\par}
\end{figure}
\end{document}
